\documentclass[10pt,letterpaper]{article}
\usepackage{cornell-notes}

\title{Clause 23.03.04: Non Member Comparison Functions}
\author{}
\date{}
\setDocTitle{Clause 23.03.04: Non Member Comparison Functions}
\setDocSubtitle{C++ 2024}
\setDocOwner{C++ 2024 Cornell Notes}
\setDocDate{\today}
\setCornellCollection{C++ 2024 Cornell Notes}
\setCornellUnitType{Clause}
\setCornellUnitNumber{23.03.04}
\setCornellUnitTitle{Non Member Comparison Functions}
\hypersetup{pdftitle={Clause 23.03.04: Non Member Comparison Functions},pdfsubject={Cornell notes},pdfauthor={}}

\begin{document}
\maketitle
\begin{CornellOverviewBox}{Overview}
\textbf{Classification:} Standard library - Normative\\[2pt]
\textbf{Learning objective:} Understand the rules, terminology, and semantic consequences associated with Non-member comparison functions.
\end{CornellOverviewBox}\begin{CornellSummaryBox}{Summary}
{\sffamily\bfseries\color{Accent} Summary}\par\vspace{3pt}Section 23.3.4 focuses on Non-member comparison functions. Apply the specified interface together with its mandates, effects, result conditions, exception behavior, and complexity constraints. When applying it, preserve the distinction between mandatory requirements, recommendations, permissions, and behavior deliberately left undefined, unspecified, or implementation-defined.
\end{CornellSummaryBox}

\end{document}
